\documentclass[10pt]{article}
\usepackage[T1]{fontenc}
\usepackage[utf8]{inputenc}
\usepackage[french]{babel}
\usepackage{amsmath,amssymb,mathrsfs,mathpazo}
\usepackage{geometry}
\geometry{a4paper,margin=2cm}
\usepackage{graphicx,array,tabularx,tabularray,colortbl,multirow,multicol,booktabs,enumitem,pifont,ulem,fancybox,diagbox,cancel}
\usepackage{pstricks,pstricks-add,pst-plot,pst-tree,pst-node,pst-text,pst-func,pst-math,pst-bspline,pst-3dplot,pst-eucl,pst-solides3d}
\let\pscalculate\undefined
\usepackage{pst-bezier}
\usepackage{tikz,pgfplots,tkz-tab}
\usetikzlibrary{arrows,arrows.meta,patterns,calc,positioning,trees,decorations.markings,shapes.geometric,shapes.misc,shapes.arrows,intersections,decorations.pathmorphing,decorations.pathreplacing,backgrounds,shadows,babel}
\pgfplotsset{compat=1.18}
\usepgfplotslibrary{fillbetween}
\usepackage[np]{numprint}
\usepackage{listings,makecell,pas-tableur,twoopt}
\let\barre\undefined
\usepackage[pstricks]{tablvar}
\definecolor{darkgreen}{RGB}{0,100,0}
\providecommand{\red}{\color{red}}
\providecommand{\blue}{\color{blue}}
\providecommand{\green}{\color{green}}
\providecommand{\black}{\color{black}}
\providecommand{\white}{\color{white}}
\providecommand{\gray}{\color{gray}}
\providecommand{\cyan}{\color{cyan}}
\providecommand{\magenta}{\color{magenta}}
\providecommand{\yellow}{\color{yellow}}
\providecommand{\violet}{\color{violet}}
\providecommand{\orange}{\color{orange}}
\providecommand{\euro}{\texteuro}
\providecommand{\EUR}{\texteuro}
\providecommand{\ieme}{\textsuperscript{e}}
\providecommand{\degres}{\ensuremath{^\circ}}
\providecommand{\degree}{\ensuremath{^\circ}}
\providecommand{\vv}[1]{\overrightarrow{#1}}
\providecommand{\Coord}[1]{#1}
\providecommand{\decosix}{}
\providecommand{\starredbullet}{\ensuremath{\bullet}}
\pagestyle{empty}
\setlength{\parindent}{0pt}
\usepackage{adjustbox}\DeclareUnicodeCharacter{2236}{:}\DeclareUnicodeCharacter{2190}{\ensuremath{\leftarrow}}
\providecommand{\R}{\mathbb{R}}
\providecommand{\N}{\mathbb{N}}
\providecommand{\Z}{\mathbb{Z}}
\providecommand{\Q}{\mathbb{Q}}
\providecommand{\D}{\mathbb{D}}
\providecommand{\C}{\mathbb{C}}
\providecommand{\e}{{\mathrm{e}}}
\providecommand{\vect}[1]{\overrightarrow{#1}}
\providecommand{\vectt}[1]{\overrightarrow{\mathrm{#1}}}
\providecommand{\vvt}[1]{\overrightarrow{\mathrm{#1}}}
\providecommand{\np}[2][]{{#2~#1}}
\providecommand{\vv}[1]{{\overrightarrow{#1}}}
\renewcommand{\d}{{\mathrm{d}}}
\providecommand{\degres}{{^\circ}}
\providecommand{\degre}{{^\circ}}
\providecommand{\blue}{\color{blue}}
\providecommand{\red}{\color{red}}
\providecommand{\enskip}{\quad}
\providecommand{\llbracket}{\lbrack\!\lbrack}
\providecommand{\rrbracket}{\rbrack\!\rbrack}
\providecommand{\euro}{{\text{€}}}
\providecommand{\no}{{\text{n°}}}
\providecommand{\interval}[2]{{[#1;#2]}}
\providecommand{\intervFO}[2]{{[#1;#2[}}
\providecommand{\dotfill}{\ldots}

\providecommand{\ds}{\displaystyle}
\providecommand{\textding}[1]{#1}
\providecommand{\Oij}{\(\mathrm{O};\vec{\imath},\vec{\jmath}\)}
\providecommand{\Oijk}{\(\mathrm{O};\vec{\imath},\vec{\jmath},\vec{k}\)}
\providecommand{\Ouv}{\(\mathrm{O};\vec{u},\vec{v}\)}
\providecommand{\intervFF}[2]{\ensuremath{[#1;#2]}}
\renewcommand{\euro}{\ifmmode\text{\texteuro}\else\texteuro\fi}
\DeclareUnicodeCharacter{00B0}{\ifmmode^{\circ}\else\textdegree\fi}
\DeclareUnicodeCharacter{20AC}{\ifmmode\text{\texteuro}\else\texteuro\fi}
% Énoncés : Éducation nationale. Transcription et corrigés : APMEP, Denis Vergès et contributeurs. Découpage et classement : CoopMaths, Sébastien Lozano, Éric Elter et contributeurs. Adaptation de la présentation : Maths Explorer.
% CC BY-SA : https://coopmaths.fr/www/about/
% https://coopmaths.fr/alea/static/eam/2026/tex/eam_specifique_2026_06_etrangers_automatismes_cor.tex
\begin{document}
\noindent\textit{APMEP / CoopMaths. CC BY-SA. Présentation adaptée par Maths Explorer.}\par\medskip
\textbf{Pour cette première partie, aucune justification n'est demandée et une seule réponse est possible par question. Pour chaque question, reporter son numéro sur la copie et indiquer la réponse.}
	
	Une réponse fausse ou l'absence de réponse n'enlève aucun point.
	
	\medskip
	
	\textbf{Question 1}\quad On considère $A=4-2\times\dfrac{1}{3}$. On a :
	
	\begin{tblr}{XXXX}
		\textbf{a.} $A=\dfrac{2}{3}$ & \textbf{b.} $A=\dfrac{10}{3}$ & \textbf{c.} $A=\dfrac{4}{3}$ & \textbf{d.} $A=\dfrac{11}{3}$ \\
	\end{tblr}
	
	\begin{color}{blue}
	\textbf{b.} $A=\dfrac{10}{3}$
	 
	 $A=4-2\times\dfrac{1}{3} = \dfrac{12}{3} - \dfrac{2}{3} = \dfrac{10}{3}$
	\end{color}
	
	\medskip
	
	\textbf{Question 2}\quad On considère $B=2\times 5^{2}+3$. On a :
	
	\begin{tblr}{XXXX}
		\textbf{a.} $B=103$ & \textbf{b.} $B=53$ & \textbf{c.} $B=97$ & \textbf{d.} $B=23$ \\
	\end{tblr}
	
	\begin{color}{blue}
	\textbf{b.} $B=53$
	 
	 $B=2\times 5^{2}+3 = 2\times 25 + 3 = 50+3 = 53$
	\end{color}
	
	
	\medskip
	
	\textbf{Question 3}\quad $25~\%$ de 250 est égal à :
	
	\begin{tblr}{XXXX}
		\textbf{a.} $62{,}5$ & \textbf{b.} 125 & \textbf{c.} 50 & \textbf{d.} 225 \\
	\end{tblr}
	
	\begin{color}{blue}
	\textbf{a.} $62{,}5$
	 
	 $25~\%\times 250 = \dfrac{1}{4}\times 250 = \np{62.5}$
	\end{color}
	
	\medskip
	
	\textbf{Question 4}\quad Un article coûtant 300~\euro{} subit une baisse de $15~\%$.
	
	Pour obtenir le prix de cet article après la baisse, il faut faire le calcul :
	
	\begin{tblr}{XXXX}
		\textbf{a.} $300-0{,}15$ & \textbf{b.} $300\times 0{,}85$ & \textbf{c.} $300\times 1{,}15$ & \textbf{d.} $300\times 0{,}15$ \\
	\end{tblr}
	
	\begin{color}{blue}
	\textbf{b.} $300\times 0{,}85$
	 
	 Le coefficient multiplicateur correspondant à une baisse de 15~\% est donné par $1-15~\% = 1-\np{0.15} = \np{0.85}$
	\end{color}
	
	\medskip
	
	\begin{minipage}{0.58\linewidth}
		\textbf{Question 5}
		
		Dans un repère du plan, on a représenté la droite (AB).
		
		L'équation réduite de la droite (AB) est :
		
		\begin{tblr}{XX}
			\textbf{a.} $y=4x+2$ & \textbf{b.} $y=-2x+2$ \\
			\textbf{c.} $y=2x+4$ & \textbf{d.} $y=-0{,}5x+2$ \\
		\end{tblr}
	\end{minipage}
	\hfill
	\begin{minipage}{0.38\linewidth}
		\begin{center}
			\begin{tikzpicture}[scale=0.5]
				\draw[gray!60,step=1,dashed] (-2,-1.7) grid (6.3,4.5);
				\draw[-latex] (-2,0)--(6.3,0);
				\draw[-latex] (0,-1.7)--(0,4.5);
				\foreach \x in {-2,-1,1,2,3,4,5,6} \draw (\x,2pt)--(\x,-2pt) node[below]{\scriptsize $\x$};
				\foreach \y in {-1,1,2,3,4} \draw (2pt,\y)--(-2pt,\y) node[left]{\scriptsize $\y$};
				\node[below left] at (0,0) {\scriptsize 0};
				\draw[blue,line width=1pt] (-2,3)--(6,-1);
				\fill (0,2) circle (2.5pt); \node[above right] at (0,2) {\scriptsize A};
				\fill (4,0) circle (2.5pt);   \node[above right] at (4,0) {\scriptsize B};
			\end{tikzpicture}
		\end{center}
	\end{minipage}
	
	\begin{color}{blue}
	\textbf{d.} $y=-0{,}5x+2$
	 
	 Par lecture des coordonnées de A$(0,\mathbf{2})$ on peut donner l'ordonnée à l'origine $2$. Le coefficient directeur s'obtient par le calcul : 
	 
	 $ \dfrac{0 - 2}{4 - 0} = \dfrac{-2}{4} = -\dfrac{1}{2} = \np{-0.5} $
	 
	 Ainsi $y=-0{,}5x+2$
	\end{color}
	
	\medskip
	
	\textbf{Question 6}\quad La valeur de l'expression \quad $2x^{2}-3x-4$ \quad pour $x=-1$ est :
	
	\begin{tblr}{XXXX}
		\textbf{a.} $-9$ & \textbf{b.} $-3$ & \textbf{c.} $-5$ & \textbf{d.} 1 \\
	\end{tblr}
	
	\begin{color}{blue}
	\textbf{d.} 1
	
	$2\times(-1)^{2} - 3\times(-1) - 4 = 2 + 3 - 4 = 1$
	\end{color}
	
	\medskip
	
	\textbf{Question 7}\quad L'expression \quad $(x-4)^{2}$ \quad est égale à :

	\begin{tblr}{XX}
		\textbf{a.} $x^{2}-8x+16$ & \textbf{b.} $x^{2}+8x+16$ \\
		\textbf{c.} $x^{2}-8x-16$ & \textbf{d.} $x^{2}+8x-16$ \\
	\end{tblr}
	
	\begin{color}{blue}
	\textbf{a.} $x^{2}-8x+16$
	
	$ (x-4)^{2} = x^{2} - 2\times x \times 4 + 4^2 = x^{2}-8x+16 $
	\end{color}
	
	\medskip
	
	\begin{minipage}{0.5\linewidth}
		\textbf{Question 8}
		
		Voici la courbe représentative d'une fonction $f$ définie sur $[-6\,;\,5]$.
		
		L'ensemble $\mathcal{S}$ des solutions de l'inéquation \quad $f(x)\geqslant 3$ \quad est :
		
		\begin{enumerate}[label=\textbf{\alph*.},leftmargin=*]
			\item $\mathcal{S}=[-6\,;-5]\cup[-2\,;5]$
			\item $\mathcal{S}=\{-5\,;-2\}$
			\item $\mathcal{S}=[-5\,;-2]$
			\item $\mathcal{S}=\{-3\}$
		\end{enumerate}
	\end{minipage}
	\hfill
	\begin{minipage}{0.46\linewidth}
		\begin{center}
			\begin{tikzpicture}
				\begin{axis}[/pgf/number format/.cd, use comma,
					width=7.5cm,height=6cm,
					xmin=-6.5,xmax=5.5,ymin=-4,ymax=5,
					xtick={-6,-5,-4,-3,-2,-1,1,2,3,4,5},
					ytick={-3,-2,-1,1,2,3,4},
					tick label style={font=\scriptsize},
					grid=both, axis lines=center,
					grid style={gray!60,line width=0.3pt,dashed}]
					\addplot[line width=1.1pt,color=blue,smooth,samples=200,domain=-6:5]
					{(0.05*x^3+0.05*x^2-1.6*x)};
					\addplot[line width=2.5pt,color=green,opacity=0.5,smooth,samples=200,domain=-5:-2]
					{(0.05*x^3+0.05*x^2-1.6*x)};
					\draw[line width=3.5pt,color=green,opacity=0.5] (-5,0) -- (-2,0);
					\draw[line width=1.5pt,color=teal,opacity=0.7,dotted] (-5,0) -- (-5,3) (-2,3) -- (-2,0) (-6,3) -- (5,3);
				\end{axis}
			\end{tikzpicture}
		\end{center}
	\end{minipage}
	
	
	\begin{color}{blue}
	\textbf{c.} $\mathcal{S}=[-5~;~-2]$
	\end{color}
	
	\medskip
	
	
	\textbf{Question 9}\quad L'ensemble $\mathcal{S}$ des solutions de l'équation \quad $(2x+4)(-3x-9)=0$ \quad est :
	
	\begin{tblr}{XXXX}
		\textbf{a.} $\mathcal{S}=\{-5\}$ & \textbf{b.} $\mathcal{S}=\{-4\,;9\}$ & \textbf{c.} $\mathcal{S}=\{-2\,;3\}$ & \textbf{d.} $\mathcal{S}=\{-3\,;-2\}$ \\
	\end{tblr}
	
	\begin{color}{blue}
	\(
	\begin{array}{ccc}
	(2x+4)(-3x-9)=0 &\iff& (2x+4)=0\text{ ou }(-3x-9)=0 \\
	(2x+4)(-3x-9)=0 &\iff& 2x=-4\text{ ou }-9=3x \\
	(2x+4)(-3x-9)=0 &\iff& x=-2\text{ ou }-3=x \\
	\end{array}
	\)
	
	\textbf{d.} $\mathcal{S}=\{-3\,;-2\}$
	
	\end{color}
	
	\medskip
	
	\textbf{Question 10}\quad Soit la formule \quad $F=G\times\dfrac{m_{1}\times m_{2}}{R^{2}}$.\quad On a :
	

	\begin{tblr}{XX}
		\textbf{a.} $m_{1}=\dfrac{F\times G}{R^{2}\times m_{2}}$ & \textbf{b.} $m_{1}=\dfrac{F\times R^{2}}{G\times m_{2}}$\\ 
		\textbf{c.} $m_{1}=\sqrt{\dfrac{G\times R^{2}}{F\times m_{2}}}$ & \textbf{d.} $m_{1}=F\times R^{2}\times G\times m_{2}$ \\
	\end{tblr}


	\begin{color}{blue}
En supposant que toutes les grandeurs sont non-nulles :	
	
	\(
	\begin{array}{ccc}
	F=G\times\dfrac{m_{1}\times m_{2}}{R^{2}} &\iff& F \times R^{2} = G\times m_{1}\times m_{2}\\
	&\iff& F \times R^{2} =  m_{1}\times(G \times m_{2})\\
	&\iff& \dfrac{F \times R^{2}}{G\times m_{2}} = m_{1}
	\end{array}
	\)
	
	\textbf{b.} $m_{1}=\dfrac{F\times R^{2}}{G\times m_{2}}$
	
	\end{color}



	\medskip
	
	Pour les questions 11 et 12, $A$ et $B$ sont deux événements et on considère l'arbre pondéré suivant :
	
	\begin{center}
		\begin{tikzpicture}[xscale=1,yscale=1]
			\tikzstyle{noeud}=[fill=white,inner sep=1pt]
			\tikzstyle{etiquette}=[pos=0.5,fill=white, inner xsep=2pt, inner ysep=1pt]
			\tikzstyle{etiquettesol}=[pos=0.5,preaction={fill=blue!20!white,opacity=0.8},color =blue,circle, inner xsep=2pt, inner ysep=1pt]
			\def\NA{0}\def\NB{2.4}\def\NC{4.8}
			\node[noeud] (R)  at (\NA,0) {};
			\node[noeud] (A)  at (\NB,1.4) {$A$};
			\node[noeud] (Ab) at (\NB,-1.4) {$\overline{A}$};
			\node[noeud] (B1) at (\NC,2.2) {$B$};
			\node[noeud] (Bb1)at (\NC,0.6) {$\overline{B}$};
			\node[noeud] (B2) at (\NC,-0.6) {$B$};
			\node[noeud] (Bb2)at (\NC,-2.2) {$\overline{B}$};
			\draw (R)--(A) node[etiquettesol]{$0{,}2$}; 
			\draw (R)--(Ab)  node[etiquette,]{$0{,}8$};
			\draw (A)--(B1)  node[etiquette]{$0{,}3$};
			\draw (A)--(Bb1) node[etiquettesol,draw=blue,line width=1.5pt]{$0{,}7$};
			\draw (Ab)--(B2) node[etiquettesol]{$0{,}4$};
			\draw (Ab)--(Bb2)node[etiquette,draw=blue,circle,draw,line width=1.5pt]{$0{,}6$};
		\end{tikzpicture}
	\end{center}
	
	\textbf{Question 11}\quad On a :
	
	\begin{tblr}{XXXX}
		\textbf{a.} $\text{P}_{\overline{A}}(\overline{B})=0{,}3$ & \textbf{b.} $\text{P}_{\overline{A}}(\overline{B})=0{,}48$ & \textbf{c.} $\text{P}_{\overline{A}}(\overline{B})=0{,}8$ & \textbf{d.} $\text{P}_{\overline{A}}(\overline{B})=0{,}6$ \\
	\end{tblr}
	
	\begin{color}{blue}
	\textbf{d.} $\text{P}_{\overline{A}}(\overline{B})=0{,}6$
	
	Probabilité de $\overline{B}$ sachant $\overline{A}$.
	
	\end{color}
	
	\medskip
	
	\textbf{Question 12}\quad On a :
	
	\begin{tblr}{XXXX}
		\textbf{a.} $\text{P}_{A}(\overline{B})=0{,}3$ & \textbf{b.} $\text{P}_{A}(\overline{B})=0{,}2$ & \textbf{c.} $\text{P}_{A}(\overline{B})=0{,}14$ & \textbf{d.} $\text{P}_{A}(\overline{B})=0{,}7$ \\
	\end{tblr}
	
	\begin{color}{blue}
	\textbf{d.} $\text{P}_{A}(\overline{B})=0{,}7$
	
	Probabilité de $\overline{B}$ sachant ${A}$.
	\end{color}

\end{document}
